\documentclass[11pt]{article}

\usepackage[margin=1in]{geometry}
\usepackage{amsmath}
\usepackage{amssymb}
\usepackage{bm}
\usepackage{hyperref}
\usepackage{xcolor}
\usepackage{listings}

\hypersetup{
  colorlinks=true,
  linkcolor=blue,
  urlcolor=blue,
  citecolor=blue
}

\lstset{
  basicstyle=\ttfamily\small,
  breaklines=true,
  frame=single,
  columns=fullflexible,
  keepspaces=true
}

\newcommand{\dd}{\mathrm{d}}
\newcommand{\ub}{\mu\mathrm{b}}
\newcommand{\MeV}{\mathrm{MeV}}
\newcommand{\GeV}{\mathrm{GeV}}
\newcommand{\sr}{\mathrm{sr}}
\newcommand{\eps}{\varepsilon}
\newcommand{\phicm}{\phi_{\mathrm{cm}}}
\newcommand{\thecm}{\theta_{\mathrm{cm}}}
\newcommand{\xb}{x_B}
\newcommand{\Qsq}{Q^2}
\newcommand{\Ebeam}{E}
\newcommand{\Escat}{E'}
\newcommand{\Mtar}{M}

\title{SIMC Cross-Section and Normalization Notes for Exclusive, SIDIS, and Delta-Pi0 Modes}
\author{NPS analysis notes}
\date{\today}

\begin{document}
\maketitle

\section{Purpose}

This note summarizes the SIMC normalization and cross-section flow discussed for
the NPS pi0 analysis, with particular emphasis on the relation between the
exclusive SIMC 5-fold cross section
\[
  \frac{\dd^5\sigma}{\dd \Escat\,\dd\Omega_e\,\dd\Omega_h}
\]
and the differential form
\[
  \frac{\dd\sigma}
       {\dd \Qsq\,\dd \xb\,\dd\phi_e\,\dd t\,\dd\phicm}.
\]

The code references below refer to
\texttt{/group/nps/singhav/simc\_gfortran\_updated}.

\section{Main SIMC Weighting Flow}

SIMC generates Monte Carlo events in spectrometer variables, computes a physics
cross section at the event kinematics, and then scales the accepted event sum by
luminosity and generated phase-space volume.

At the event level, the physics cross section is stored in
\texttt{main\%sigcc}.  The dispatch is in \texttt{event.f}:

\begin{lstlisting}
1458 elseif (doing_pion) then
1459   main%sigcc = peepi(vertex,main)
...
1511 elseif (doing_delta) then
1512   main%sigcc = peedelta(vertex,main) !Need new xsec model.
...
1520 elseif (doing_semi) then
1521   main%sigcc = peepiX(vertex,vertex0,main,survivalprob,.FALSE.)
\end{lstlisting}

The event weight is then

\begin{lstlisting}
1558 main%weight = main%SF_weight*main%jacobian*main%gen_weight*main%sigcc
1559 main%weight = main%weight * tgtweight
\end{lstlisting}

Thus, schematically,
\[
  w_i
  =
  W_{\mathrm{SF},i}\,
  J_{\mathrm{gen},i}\,
  W_{\mathrm{gen},i}\,
  \sigma_i\,
  W_{\mathrm{target}},
\]
where
\begin{itemize}
  \item \(W_{\mathrm{SF}}\) is the spectral-function or nuclear model weight;
  \item \(J_{\mathrm{gen}}\equiv\texttt{main\%jacobian}\) converts generated
        spectrometer slopes to solid angle;
  \item \(W_{\mathrm{gen}}\equiv\texttt{main\%gen\_weight}\) accounts for event-by-event
        restrictions of the generated energy interval;
  \item \(\sigma_i\equiv\texttt{main\%sigcc}\) is the physics cross section;
  \item \(W_{\mathrm{target}}\equiv\texttt{tgtweight}\) counts the relevant target
        protons or neutrons for pion/kaon production.
\end{itemize}

\section{Luminosity, Target Factor, and Normalization}

\subsection{Target factor}

SIMC defines
\begin{lstlisting}
94 targetfac=targ%mass_amu/3.75914d+6/(targ%abundancy/100.)
95      >    *abs(cos(targ%angle))/(targ%thick*1000.)
\end{lstlisting}
in \texttt{simc.f}.  The luminosity is then
\begin{lstlisting}
101 luminosity = EXPER%charge/targetfac
\end{lstlisting}

Equivalently,
\[
  \mathcal{L}_{\mathrm{int}}
  =
  Q_{\mathrm{mC}}\,
  3.75914\times 10^6\,
  \frac{(t_{\mathrm{mg/cm^2}})\,f_{\mathrm{abund}}}
       {A_{\mathrm{amu}}\,|\cos\theta_{\mathrm{target}}|},
\]
where
\[
  f_{\mathrm{abund}} = \frac{\texttt{targ\%abundancy}}{100}.
\]
The output units are \(\ub^{-1}\).  The factor \texttt{targ\%thick*1000}
appears because the input thickness is given in \(\mathrm{mg/cm^2}\), but the
database processing converts \texttt{targ\%thick} to \(\mathrm{g/cm^2}\).

\subsection{End-of-run normalization}

The run-level normalization starts as
\begin{lstlisting}
368 normfac = luminosity/ntried*nevent
\end{lstlisting}
which Fortran evaluates left-to-right:
\[
  \texttt{normfac}
  =
  \mathcal{L}_{\mathrm{int}}\,
  \frac{N_{\mathrm{event}}}{N_{\mathrm{tried}}}.
\]
It is not \(\mathcal{L}_{\mathrm{int}}/(N_{\mathrm{tried}}N_{\mathrm{event}})\).

SIMC then multiplies by the generated phase-space volume:
\begin{lstlisting}
376 domega_e=(gen%e%yptar%max-gen%e%yptar%min)
377      >    *(gen%e%xptar%max-gen%e%xptar%min)
...
380 domega_p=(gen%p%yptar%max-gen%p%yptar%min)
381      >    *(gen%p%xptar%max-gen%p%xptar%min)
383 genvol = domega_e
...
387 if (doing_deuterium.or.doing_heavy.or.doing_pion.or.doing_kaon
388      >      .or.doing_delta.or.doing_rho .or. doing_semi) then
389   genvol = genvol * domega_p * (gen%e%E%max-gen%e%E%min)
390 endif
392 if (doing_heavy.or.doing_semi) then
393   genvol = genvol * (gen%p%E%max-gen%p%E%min)
394 endif
396 normfac = normfac * genvol
\end{lstlisting}

Therefore
\[
  \texttt{normfac}
  =
  \mathcal{L}_{\mathrm{int}}\,
  \frac{N_{\mathrm{event}}}{N_{\mathrm{tried}}}\,
  V_{\mathrm{gen}}.
\]

For exclusive pion/pi0 and the separate \texttt{doing\_delta} mode,
\[
  V_{\mathrm{gen}}^{\mathrm{excl}}
  =
  \Delta\Omega_e^{\mathrm{gen}}\,
  \Delta\Omega_h^{\mathrm{gen}}\,
  \Delta \Escat .
\]
For SIDIS,
\[
  V_{\mathrm{gen}}^{\mathrm{SIDIS}}
  =
  \Delta\Omega_e^{\mathrm{gen}}\,
  \Delta\Omega_h^{\mathrm{gen}}\,
  \Delta \Escat\,
  \Delta E_h .
\]

\section{The Three Different ``Jacobian'' Factors}

SIMC uses the name \texttt{jacobian} in several different contexts.  They are
not the same object.

\subsection{Generated slope to solid angle: \texttt{main\%jacobian}}

SIMC generates uniformly in spectrometer slopes \((x'_{\mathrm{tar}},y'_{\mathrm{tar}})\),
not directly in physical solid angle.  The code comment in \texttt{event.f} says

\begin{lstlisting}
265 ! 1: We generate uniformly in xptar/yptar, not theta/phi.
267 ! range, and have a jacobian for each event taking into account the mapping
268 ! between the solid angle on the unit sphere, and the dxptar/dyptar volume
269 ! (the jacobian is 1/cos**3(dtheta), where dtheta is the angle between the
270 ! event and the central spectrometer vector
\end{lstlisting}

The actual code is
\begin{lstlisting}
1013 r = sqrt(1.+vertex%e%yptar**2+vertex%e%xptar**2)
1014 main%jacobian = main%jacobian / r**3
...
1021 r = sqrt(1.+vertex%p%yptar**2+vertex%p%xptar**2)
1022 main%jacobian = main%jacobian / r**3
\end{lstlisting}

Thus, for each generated arm,
\[
  J_{\mathrm{gen}}
  =
  \frac{\dd\Omega}{\dd x'_{\mathrm{tar}}\,\dd y'_{\mathrm{tar}}}
  =
  \frac{1}
       {\left(1+x_{\mathrm{tar}}^{\prime\,2}
                +y_{\mathrm{tar}}^{\prime\,2}\right)^{3/2}}.
\]
This is the factor appearing in the event weight as \texttt{main\%jacobian}.

\subsection{Exclusive CM-to-lab Jacobian}

The exclusive pion model gives a hadronic virtual-photon cross section in
\((t,\phicm)\).  SIMC needs the lab hadron solid angle for its generated
variables, so \texttt{transform\_to\_cm} returns a local variable also named
\texttt{jacobian}.

The exact definition in \texttt{jacobians.f} is
\begin{lstlisting}
271 jacobian = abs(dt_dcos_lab*dphicmdphi-dt_dphi_lab*dphicmdcos)
\end{lstlisting}

Mathematically,
\[
  J_{\mathrm{excl}}
  =
  \left|
  \frac{\partial(t,\phicm)}
       {\partial(\cos\theta_{h,\mathrm{lab}},\phi_{h,\mathrm{lab}})}
  \right|.
\]
This is the factor that converts
\[
  \frac{\dd^2\sigma_{\gamma^*N}}{\dd t\,\dd\phicm}
  \quad\longrightarrow\quad
  \frac{\dd^2\sigma_{\gamma^*N}}{\dd\Omega_{h,\mathrm{lab}}}.
\]

\subsection{SIDIS \((z,p_T^2)\) to detected-hadron variables}

In \texttt{semi\_physics.f}, the SIDIS model first constructs
\[
  \frac{\dd\sigma}
       {\dd\Omega_e\,\dd \Escat\,\dd z\,\dd p_T^2\,\dd\phi_h}.
\]
It is then converted to
\[
  \frac{\dd\sigma}
       {\dd\Omega_e\,\dd \Escat\,\dd E_h\,\dd\cos\theta_{hq}\,\dd\phi_h}.
\]
The code says
\begin{lstlisting}
556 C Need to convert to dsig/ (dOmega_e dE_e dE_h dCos(theta) dPhi_had
557 C This is just given by 1/omega * 2*p_h**2*cos(theta)
559 jacobian = 1./(nu/1000.)*2.*(phad/1000.)**2*cthpq
562 sigma_eepiX = sigsemi*jacobian/1.e6
\end{lstlisting}
Thus
\[
  J_{\mathrm{SIDIS}}
  =
  \frac{1}{\nu}\,2p_h^2\cos\theta_{hq},
\]
with \(\nu\) and \(p_h\) expressed in \(\GeV\) in that line.

\section{Exclusive Pion/Pi0 Cross Section}

\subsection{Function and units}

The exclusive pion/pi0 cross section is calculated in \texttt{peepi} in
\texttt{physics\_pion.f}.  The function header states
\begin{lstlisting}
8 *   output:
9 * peepi !d5sigma/dEe'dOmegae'Omegapi (microbarn/MeV/sr^2)
\end{lstlisting}

Inside the routine, the model first produces a 2-fold virtual-photon
center-of-mass cross section:
\begin{lstlisting}
169 * sigma_eepi is two-fold C.M. cross section: d2sigma/dt/dphi_cm [ub/MeV**2/rad]
\end{lstlisting}
That quantity is called \(\sigma_{\mathrm{cm}}\) below:
\[
  \sigma_{\mathrm{cm}}
  \equiv
  \frac{\dd^2\sigma_{\gamma^*N}}
       {\dd t\,\dd\phicm}.
\]

\subsection{The active model}

For the high-\(W\) region, SIMC uses Peter Bosted's 2021 exclusive fit:
\begin{lstlisting}
125 CDG Use Peter Bosted's new fit to world, JLab 6 GeV, and preliminary 12 GeV data
126 ntup%sigcm1 = sig_param_2021(thetacm,phicm,main%t/1.e6,vertex%q2/1.e6,s/1.e6,main%epsilon,
127      >           which_pion,doing_pizero)
129 sigma_eepi = ntup%sigcm1
\end{lstlisting}

For \(\pi^0\), \texttt{sig\_param\_2021} averages the pi-plus and pi-minus fits:
\begin{lstlisting}
788 if(doing_pizero) then
789    call exclfit(... sigexcl3p, pp, 0.0)
791    call exclfit(... sigexcl3m, pm, 0.0)
793    sigexcl3 = (sigexcl3m + sigexcl3p)/2.0 ! first pass, average of pi+ and pi-
\end{lstlisting}

Below \(W=2~\GeV\), SIMC blends toward the \texttt{sigmaid} MAID-style model:
\begin{lstlisting}
133 if(main%wcm.lt.2000) then ! W less than 2.0 GeV
146    call sigmaid(ipi,Q2gev,Wgev,E0,cthcm,phicm,sig0)
148    ntup%sigcm2 = sig0 / ppicm / qstar / 2.
152    fac1 = min(1., max(0.,(Wgev - 1.5) / 0.4))
153    sigma_eepi = ntup%sigcm1 * fac1 + ntup%sigcm2 * (1 - fac1)
154 endif
\end{lstlisting}

Thus
\[
  \sigma_{\mathrm{cm}}
  =
  \begin{cases}
    \sigma_{\mathrm{MAID}}, & W \lesssim 1.5~\GeV,\\
    f(W)\,\sigma_{\mathrm{Bosted2021}}
      +[1-f(W)]\,\sigma_{\mathrm{MAID}},
      & 1.5~\GeV < W < 1.9~\GeV,\\
    \sigma_{\mathrm{Bosted2021}}, & W \gtrsim 1.9~\GeV,
  \end{cases}
\]
where
\[
  f(W)=\min\left[1,\max\left(0,\frac{W-1.5}{0.4}\right)\right].
\]

\subsection{Separated 2-fold form}

The Bosted 2021 fit is evaluated through \texttt{exclfit}.  Its separated
structure is
\begin{lstlisting}
840 sig219 = (sigt +
841      >   eps * sigl +
842      >   eps * cos(2.0*phicm) * sigtt +
843      >   sqrt(2.0 * eps * (1.0+eps)) *
844      >   cos(phicm) * siglt)
850 sig = sig219 *8.539 / (s_gev - mtar_gev**2)**2
851 sig = sig / 2.0 / 3.1415928 / 1.0d+06
\end{lstlisting}

In mathematical notation,
\[
  \frac{\dd^2\sigma_{\gamma^*N}}
       {\dd t\,\dd\phicm}
  =
  \left[
    \sigma_T
    +\eps\,\sigma_L
    +\eps\cos(2\phicm)\,\sigma_{TT}
    +\sqrt{2\eps(1+\eps)}\cos\phicm\,\sigma_{LT}
  \right]
  \frac{8.539}{(W^2-M_p^2)^2}
  \frac{1}{2\pi}
  \frac{1}{10^6}.
\]
The final factor \(10^{-6}\) converts the result to \(\ub/\MeV^2/\mathrm{rad}\)
when \(t,Q^2,W^2\) are passed in \(\GeV\)-based units inside the fit.

\subsection{Conversion to the SIMC exclusive 5-fold cross section}

The exclusive 5-fold result returned by \texttt{peepi} is
\begin{lstlisting}
188 fac = 1./(1.-pferz*pfer/efer) * targ%Mtar_struck/efer
189 gtpr = alpha/2./(pi**2)*vertex%e%E/vertex%Ein*k_eq/vertex%q2/(1.-main%epsilon)
191 peepi = sigma_eepi*jacobian*(gtpr*fac)
\end{lstlisting}

Define
\[
  \Gamma_T'
  =
  \frac{\alpha}{2\pi^2}
  \frac{\Escat}{\Ebeam}
  \frac{K_{\mathrm{eq}}}{\Qsq}
  \frac{1}{1-\eps},
\]
where the SIMC equivalent photon energy is
\[
  K_{\mathrm{eq}}
  =
  \frac{W^2-M_{\mathrm{struck}}^2}{2M_{\mathrm{struck}}}.
\]
Also define the Fermi-motion flux correction
\[
  F_{\mathrm{fermi}}
  =
  \frac{1}{1-p_{\parallel}/E_{\mathrm{fermi}}}
  \frac{M_{\mathrm{struck}}}{E_{\mathrm{fermi}}},
\]
which is implemented as the variable \texttt{fac}.  Then SIMC returns
\[
  \boxed{
  \frac{\dd^5\sigma_{\mathrm{excl}}}
       {\dd \Escat\,\dd\Omega_e\,\dd\Omega_{h,\mathrm{lab}}}
  =
  \frac{\dd^2\sigma_{\gamma^*N}}
       {\dd t\,\dd\phicm}
  \,
  J_{\mathrm{excl}}
  \,
  \Gamma_T'
  \,
  F_{\mathrm{fermi}}
  }.
\]

\section{Delta-Pi0 Channels}

There are two distinct ``delta'' meanings in this SIMC checkout.

\subsection{The NPS delta-pi0 input path}

The file \texttt{nps\_delta\_pi0\_x60\_4b.inp} uses
\[
  \texttt{doing\_pion}=1,\qquad
  \texttt{which\_pion}=2,\qquad
  \texttt{doing\_pizero}=1,\qquad
  \texttt{doing\_semi}=0.
\]
Therefore it follows the exclusive pion/pi0 function \texttt{peepi}, not the
separate \texttt{peedelta} function.

After \texttt{peepi}, SIMC applies an empirical Delta-channel scale factor:
\begin{lstlisting}
1464 if(which_pion.eq.2) then ! g* p -> pi+ Delta0, g* n -> pi+ Delta-, or g* p -> pi0 Delta+
1465    if(doing_hydpi) then
1466       if (doing_pizero) then
1467          main%sigcc = 0.55*main%sigcc ! g* p -> pi0 Delta+
\end{lstlisting}
Thus for the hydrogen \(\pi^0\Delta^+\) setting,
\[
  \boxed{
  \dd^5\sigma(p(e,e'\pi^0)\Delta^+)
  =
  0.55\,
  \dd^5\sigma_{\mathrm{excl}\,\pi^0}
  }.
\]

\subsection{The separate \texttt{doing\_delta} path}

The separate \texttt{doing\_delta} path calls \texttt{peedelta}.  The intended
form mirrors the exclusive pion form, but the active code currently ignores the
computed \(\sigma_{\mathrm{cm}}\):
\begin{lstlisting}
99  ntup%sigcm1 = sig_blok(thetacm,phicm,main%t/1.d6,vertex%q2/1.d6,s/1.d6,main%epsilon,
100      >    targ%Mtar_struck/1000.,which_pion)
101 sigma_eepi = ntup%sigcm1
...
130 ! peedelta = sigma_eepi*jacobian*(gtpr*fac)
132 peedelta = 1.0*jacobian*(gtpr*fac)
\end{lstlisting}
So the active \texttt{doing\_delta} return is
\[
  \dd^5\sigma_{\texttt{doing\_delta}}
  =
  1\cdot J_{\mathrm{excl}}\Gamma_T'F_{\mathrm{fermi}},
\]
not
\[
  \sigma_{\mathrm{cm}}J_{\mathrm{excl}}\Gamma_T'F_{\mathrm{fermi}}.
\]
The code comment in \texttt{event.f} marks this path with
\texttt{!Need new xsec model.}

\section{SIDIS Pi0 Cross Section}

The SIDIS function is \texttt{peepiX} in \texttt{semi\_physics.f}.  The header
states
\begin{lstlisting}
6 * This routine calculates p(e,e'pi+)X semi-iinclusive cross sections from
7 * the CTEQ5M parton distribution functions and a simple parameterization
8 * of the favored and unfavored fragmentation functions.
\end{lstlisting}

The active pion fragmentation uses Peter Bosted's 2021 parameterization:
\begin{lstlisting}
450 C Peter Bosted's parameterization
451 c uses a modified scaling variable (replaces z) for the fragmentation function
452 c (see Accardi et al https://arxiv.org/abs/0907.2395)
453 c new PB fit using zp for pions. This is z * D
\end{lstlisting}
For \(\pi^0\), it averages the favored and unfavored pion fragmentation pieces:
\begin{lstlisting}
480 if(doing_pizero) then ! just get averge of pi+ and pi-
481    u1 = (yf + yu) / 2.
482    d1 = (yf + yu) / 2.
483    ub = (yf + yu) / 2.
484    d1 = (yf + yu) / 2.
485 endif
\end{lstlisting}

The quark-weighted fragmentation ratio is
\begin{lstlisting}
499 dsigdz = (qu**2 * u    * u1 +
500      >            qu**2 * ubar * ub +
501      >      qd**2 * d    * d1 +
502      >            qd**2 * dbar * db +
503      >      qs**2 * s    * s1 +
504      >            qs**2 * sbar * sb)/sum_sq/zhad
\end{lstlisting}
and the transverse-momentum dependence is
\begin{lstlisting}
516 b =  1./ (0.120 * zhad**2 + 0.200)
517 sighad = dsigdz*b*exp(-b*pt2gev)/2./pi
\end{lstlisting}

The inclusive electron part uses \texttt{F1F2IN21}:
\begin{lstlisting}
537 call F1F2IN21(z8, a8, Q2gev, Wsq, F1, F2)
...
550 sige = 4.*alpha**2*(Eprime/1000)**2/Q2gev**2 * ( W2*W2coeff + 2.*W1*sin2th2)
\end{lstlisting}

SIMC then constructs
\begin{lstlisting}
552 C DJG This dsig/(dOmega_e dE_e dz dpt**2 dPhi_had) in microbarn/GeV**3/sr
553 sigsemi = sige*sighad*(hbarc/1000.)**2*10000.0
\end{lstlisting}
and converts to detected hadron variables:
\begin{lstlisting}
559 jacobian = 1./(nu/1000.)*2.*(phad/1000.)**2*cthpq
562 sigma_eepiX = sigsemi*jacobian/1.e6
584 sigma_eepiX = sigma_eepiX*fac
595 peepiX = sigma_eepiX
\end{lstlisting}

Thus, schematically,
\[
  \dd\sigma_{\mathrm{SIDIS}}
  =
  \left[
    \dd\sigma_e(F_1,F_2)
  \right]
  \left[
    D_{q\to h}(z,p_T^2)
  \right]
  J_{\mathrm{SIDIS}}
  F_{\mathrm{fermi}}.
\]

\section{Deriving the Desired Exclusive Form}

\subsection{Starting point}

SIMC's exclusive pion function returns
\[
  \frac{\dd^5\sigma_{\mathrm{excl}}}
       {\dd \Escat\,\dd\Omega_e\,\dd\Omega_{h,\mathrm{lab}}}
  =
  \sigma_{\mathrm{cm}}\,
  J_{\mathrm{excl}}\,
  \Gamma_T'\,
  F_{\mathrm{fermi}},
\]
with
\[
  \sigma_{\mathrm{cm}}
  =
  \frac{\dd^2\sigma_{\gamma^*N}}
       {\dd t\,\dd\phicm}.
\]

If the desired variables already include \(t\) and \(\phicm\), then the
exclusive hadron Jacobian must be removed.  Therefore
\[
  \frac{\dd\sigma}
       {\dd \Escat\,\dd\Omega_e\,\dd t\,\dd\phicm}
  =
  \sigma_{\mathrm{cm}}\,
  \Gamma_T'\,
  F_{\mathrm{fermi}}.
\]

\subsection{Electron-variable transformation}

Use
\[
  \dd\Omega_e = \dd\cos\theta_e\,\dd\phi_e
\]
and fixed beam energy \(\Ebeam\).  The electron variables satisfy
\[
  \Qsq = 2\Ebeam\Escat(1-\cos\theta_e),
  \qquad
  \nu = \Ebeam-\Escat,
  \qquad
  \xb = \frac{\Qsq}{2\Mtar\nu}.
\]
Solving for the Jacobian gives
\[
  \left|
  \frac{\partial(\Escat,\cos\theta_e)}
       {\partial(\Qsq,\xb)}
  \right|
  =
  \frac{\nu}{2\Ebeam\Escat\xb}
  =
  \frac{\Qsq}{4\Mtar\Ebeam\Escat\xb^2}.
\]

Derivation:
\[
  \frac{\partial \Qsq}{\partial \Escat}
  =
  2\Ebeam(1-\cos\theta_e)
  =
  \frac{\Qsq}{\Escat},
\]
\[
  \frac{\partial \Qsq}{\partial \cos\theta_e}
  =
  -2\Ebeam\Escat,
\]
\[
  \frac{\partial \xb}{\partial \Escat}
  =
  \frac{\xb}{\Escat}+\frac{\xb}{\nu},
  \qquad
  \frac{\partial \xb}{\partial\cos\theta_e}
  =
  -\frac{\Ebeam\Escat}{\Mtar\nu}.
\]
Thus
\[
  \left|
  \frac{\partial(\Qsq,\xb)}
       {\partial(\Escat,\cos\theta_e)}
  \right|
  =
  \frac{2\Ebeam\Escat\xb}{\nu},
\]
and the inverse is
\[
  \left|
  \frac{\partial(\Escat,\cos\theta_e)}
       {\partial(\Qsq,\xb)}
  \right|
  =
  \frac{\nu}{2\Ebeam\Escat\xb}
  =
  \frac{\Qsq}{4\Mtar\Ebeam\Escat\xb^2}.
\]

\subsection{Final result}

Keeping the electron azimuth differential, the desired exclusive form is
\[
  \boxed{
  \frac{\dd\sigma}
       {\dd \Qsq\,\dd\xb\,\dd\phi_e\,\dd t\,\dd\phicm}
  =
  \left[
    \frac{\dd^2\sigma_{\gamma^*N}}
         {\dd t\,\dd\phicm}
  \right]
  \Gamma_T'
  F_{\mathrm{fermi}}
  \frac{\Qsq}{4\Mtar\Ebeam\Escat\xb^2}
  }.
\]

Equivalently, in SIMC variable names,
\[
  \boxed{
  \frac{\dd\sigma}
       {\dd \Qsq\,\dd\xb\,\dd\phi_e\,\dd t\,\dd\phicm}
  =
  \texttt{sigma\_eepi}
  \times
  \texttt{gtpr}
  \times
  \texttt{fac}
  \times
  \frac{\texttt{vertex\%Q2}}
       {4\,M\,\texttt{vertex\%Ein}\,\texttt{vertex\%e\%E}\,\texttt{vertex\%xbj}^2}
  }.
\]

If instead the electron azimuth is integrated over and the cross section is
azimuthally symmetric in \(\phi_e\), then
\[
  \boxed{
  \frac{\dd\sigma}
       {\dd \Qsq\,\dd\xb\,\dd t\,\dd\phicm}
  =
  2\pi
  \left[
    \frac{\dd^2\sigma_{\gamma^*N}}
         {\dd t\,\dd\phicm}
  \right]
  \Gamma_T'
  F_{\mathrm{fermi}}
  \frac{\Qsq}{4\Mtar\Ebeam\Escat\xb^2}
  }.
\]

\subsection{Important practical warning}

Do not include the SIMC exclusive local \texttt{jacobian} when the target
variables are already \((t,\phicm)\).  That Jacobian is only needed for
\[
  (t,\phicm)\rightarrow \Omega_{h,\mathrm{lab}}.
\]
The final \((\Qsq,\xb,\phi_e,t,\phicm)\) formula keeps the hadron variables as
\((t,\phicm)\), so the hadron-side Jacobian cancels out.

\section{Unit Conventions}

SIMC mixes MeV and GeV conventions in a controlled way:
\begin{itemize}
  \item \texttt{vertex\%Q2}, \texttt{main\%t}, \texttt{vertex\%Ein}, and
        \texttt{vertex\%e\%E} are generally in MeV-based units.
  \item The exclusive parameterization is called with
        \(\Qsq,t,W^2\) divided by \(10^6\), i.e. in \(\GeV^2\).
  \item \texttt{sig\_param\_2021} converts the returned 2-fold cross section
        to \(\ub/\MeV^2/\mathrm{rad}\) using the \(10^{-6}\) factor.
  \item The SIMC \texttt{gtpr} uses MeV units and has dimensions \(1/\MeV\).
\end{itemize}

Therefore the boxed formula above is internally consistent if
\(\Qsq, \Ebeam, \Escat,\Mtar\) are all taken in the same MeV-based convention
used by SIMC.  If one rewrites the final result in \(\GeV^2\) for \(\Qsq\) and
\(t\), the corresponding powers of \(10^3\) or \(10^6\) must be applied
consistently.

\section{Minimal Dictionary}

\begin{center}
\begin{tabular}{ll}
\hline
SIMC name & Meaning \\
\hline
\texttt{sigma\_eepi} & \(\dd^2\sigma_{\gamma^*N}/(\dd t\,\dd\phicm)\) for exclusive pion/pi0 \\
\texttt{peepi} & exclusive \(\dd^5\sigma/(\dd \Escat\,\dd\Omega_e\,\dd\Omega_h)\) \\
\texttt{peepiX} & SIDIS detected-hadron cross section \\
\texttt{gtpr} & virtual photon flux \(\Gamma_T'\) \\
\texttt{fac} & Fermi-motion flux correction \\
\texttt{jacobian} in \texttt{peepi} & \(\partial(t,\phicm)/\partial(\cos\theta_h,\phi_h)\) \\
\texttt{main\%jacobian} & generated slope to solid-angle correction \\
\texttt{genvol} & generated phase-space volume \\
\texttt{normfac} & luminosity times accepted fraction times \texttt{genvol} \\
\hline
\end{tabular}
\end{center}

\end{document}
